% Capability Envelopes -- OSDI (USENIX) format
\documentclass[letterpaper,twocolumn,10pt]{article}
\usepackage{usenix-2020-09}

\usepackage{graphicx}
\usepackage{booktabs}
\usepackage{listings}
\usepackage{amsmath}

\title{Capability Envelopes: Deployable Budget Enforcement\\for Unmodified Linux Binaries}
% Double-blind submission: anonymize authors.  Camera-ready restores:
% \author{A20OS Research Team\\ \texttt{research@a20os.dev}}
\author{Anonymous Author(s)}
\date{}

\begin{document}
\maketitle

\begin{abstract}
Capability systems promise least-privilege execution but have seen
minimal adoption in production Linux environments because they demand
source-code modification---a ``rewrite mandate'' that is prohibitive for
legacy binaries.  We observe that while binary permissions are hard to
deploy retroactively, \emph{quantified budgets} are not: a budget can be
attached to any process at launch time without modifying its code.  We
present \emph{capability envelopes}, a kernel-level mechanism that wraps
unmodified Linux binaries in programmable budget policies governing type
membership, direction enforcement, data volume, operation count, and
temporal expiry.

We implement capability envelopes in A20OS, a research operating system
with a dual-ABI architecture that runs unmodified musl-linked binaries
alongside native capability-aware services.  Our kernel mediator enforces
budget accounting across ten escape-surface classes covering 366 registered
system calls, with per-direction right checks on tracked shadow handles.
A 20-cell pilot experiment across four defense configurations
(unrestricted, Landlock-permissive, Landlock-strict, envelope)
demonstrates the coarse-grained dilemma: permissive path-based policies
admit all tested attacks while strict variants block legitimate installs.
Within this pilot, only the envelope arm simultaneously completes the
benign installation and blocks all four tested attack classes, through
type denial, rights clamping, and budget exhaustion.  An
exact-execution experiment runs the unmodified install scripts of 66
real malicious packages on stock Node.js and CPython under A20OS:
every one of the 25 samples that attempted a network operation is
blocked at \texttt{socket()} under an install-time envelope, while the
same payloads reach live C2 infrastructure undefended; a complementary
canonical-replay study covers 87 samples with identical conclusions.
Overhead microbenchmarks under full-system emulation show $+4.1\%$ on
the acquisition path and $+17.7$--$29.2\%$ on the file use path.
\end{abstract}

%%------------------------------------------------------------------------
\section{Introduction}

Capability-based security systems---from Capsicum~\cite{watson2010capsicum}
to Landlock~\cite{landlock2023}---promise least-privilege execution by
replacing ambient authority with explicit capabilities.  Despite four
decades of research, adoption in production Linux remains negligible.
The root cause is what we call the \emph{rewrite mandate}: enjoying
capability discipline requires modifying source code to replace global
namespace accesses with capability-constrained operations.  For legacy
binaries, proprietary software, and third-party packages, this mandate
is prohibitive.

We observe an asymmetry that suggests a path forward.  Binary
permissions (``this process may open any file'') are difficult to
constrain retroactively because they are woven into the program's logic.
Quantified budgets (``this process may perform at most $N$ file
operations within $T$ seconds''), by contrast, can be attached at
deployment time as a wrapper policy---the binary itself needs no
modification.  We call this property \emph{budget deployability}.

\subsection{Capability Envelopes}

We present \emph{capability envelopes}: a kernel-level mechanism that
wraps unmodified Linux binaries in programmable budget policies
governing five dimensions:

\begin{enumerate}
\item \textbf{Type membership}---which resource classes (file, socket,
      pipe, memory) the process may access;
\item \textbf{Rights ceiling}---per-class maximum permission bits
      (read, write, stat, seek, connect, accept);
\item \textbf{Data budget}---cumulative bytes transferred through
      mediated I/O paths;
\item \textbf{Operation budget}---total count of mediated resource
      operations; and
\item \textbf{Temporal expiry}---wall-clock deadline after which all
      mediated operations fail atomically.
\end{enumerate}

Envelopes are implemented in A20OS, a research operating system with a
dual-ABI architecture: a Linux compatibility layer (361 registered
syscalls) runs unmodified musl-linked binaries alongside a native
capability-aware ABI (136 syscalls) for trusted components.  A kernel
mediator (approximately 670 lines of C) intercepts every resource acquisition
and consumption path, enforcing budget policies via per-handle shadow
entries with direction-aware right checks.

\subsection{Contributions}

\begin{itemize}
\item A dual-ABI kernel architecture that attaches quantified budget
      policies to unmodified binaries at launch time (\S\ref{sec:design});
\item An A1--A10 entry-point taxonomy covering 366 registered syscalls,
      mechanically verified against the syscall table via a drift gate
      (\S\ref{sec:impl});
\item A 15-scenario exercise suite (12 attack scenarios, 2 benign
      multi-process flows, 1 runtime audit) covering file, network,
      and IPC surfaces (\S\ref{sec:eval});
\item An exact-execution evaluation: 66 real malicious packages run
      unmodified on stock Node.js/CPython under A20OS---all 25 samples
      that reach the network at runtime are blocked at
      \texttt{socket()}, while reaching live C2 infrastructure
      undefended---plus an 87-sample canonical-replay study and 20
      benign install footprints completing under budget
      (\S\ref{sec:corpus});
\item A 20-cell pilot experiment across four defense arms quantifying
      the coarse-grained dilemma and showing envelopes resolve it
      (\S\ref{sec:eval});
\item A TLA+ model of the mediator state machine, exhaustively checked
      over a bounded state space (4{,}620 distinct states), plus core
      Lean~4 proofs covering monotone decay, transfer clamping, delegation bounding, and cross-envelope budget bounds
      (\S\ref{sec:formal}).
\end{itemize}

%%------------------------------------------------------------------------
\section{Background and Motivation}
\label{sec:bg}

\subsection{The Rewrite Mandate}

Capsicum~\cite{watson2010capsicum} introduced capability mode to
FreeBSD, requiring programs to enter a restricted mode and operate
exclusively through capability descriptors.  Adopting Capsicum demands
rewriting file access to use \texttt{openat(2)} with directory
descriptors, replacing \texttt{execve(2)} with capability-aware
variants, and restructuring privilege separation.  These changes are
practical only for programs designed with capability separation in mind.

Landlock~\cite{landlock2023} brought unprivileged access control to
Linux through stacked filesystem rules.  While Landlock does not
require source modification, it operates on \emph{paths}, not
\emph{objects}: a rule allowing writes to \texttt{/var/log} grants
access to every file beneath that prefix regardless of ownership or
type, and cannot express ``at most 100 writes'' or ``expires in 300
seconds.''

seccomp-BPF~\cite{seccomp2012} filters system calls by number and
argument predicates.  It cannot distinguish between opening
\texttt{/tmp/build/output.o} and \texttt{\~{}/.ssh/id\_rsa}---both are
\texttt{openat(2)} calls with different path arguments, which BPF
programs cannot reliably inspect due to pointer indirection.

The result is a dilemma we call the \emph{coarse-grained trap}: 
path-based and syscall-number-based policies face a binary choice
between permissiveness (which admits attacks) and restrictiveness
(which blocks legitimate workflows).  Our pilot experiments
(\S\ref{sec:eval}) quantify this dilemma empirically.

\subsection{Threat Model}

We target \emph{install-time attacks} in package managers (npm
postinstall, pip setup.py, cargo build scripts), CI/CD runners, and
plugin loading systems.  In these scenarios, untrusted code executes
with the user's full privileges during package installation.  Empirical
studies of malicious packages in npm and PyPI find that the majority
activate during the install or import phase rather than during normal
runtime execution~\cite{duan2021measuring,ohm2020backstabber,ladisa2023taxonomy}.

The attacker controls the package content but not the OS configuration.
The defender (user or CI operator) seeks to constrain the installation
process without breaking legitimate packages.  The defender can modify
the \emph{launch command} (adding envelope parameters) but cannot
modify the \emph{package source code}.

We do not address:
\begin{itemize}
\item Kernel vulnerabilities exploited to escape the mediation layer;
\item Side-channel information leakage through shared microarchitectural state;
\item Supply-chain compromise upstream of the package registry;
\item Social engineering convincing users to weaken envelope policies.
\end{itemize}

\subsection{Why Not Existing Linux Mechanisms?}
\label{sec:positioning}

A natural question is whether the same policy expressiveness could be
deployed on commodity Linux without a new kernel.  We survey the four
closest mechanisms and the structural gap each leaves open.

\paragraph{rlimits and cgroups.}
POSIX resource limits~\cite{setrlimit2017} and cgroups
v2~\cite{cgroupsv2} are the budget mechanisms Linux already deploys.
Both constrain \emph{resource consumption} (file size, memory, CPU,
I/O bandwidth) rather than \emph{authority}: an rlimit cannot express
``read but not write,'' neither mechanism reasons about object types or
descriptor transfers between processes, and neither offers temporal
expiry with atomic revocation.  Resource containers~\cite{banga1999rc}
established the value of kernel-visible resource principals; envelopes
borrow that insight but budget \emph{operations on authorities} rather
than scheduler-visible consumption.

\paragraph{eBPF-LSM.}
BPF LSM (KRSI)~\cite{krsi2019} attaches custom policy to object-level
security hooks and is the closest programmable neighbor.  Three
structural differences remain.  First, \emph{deployment privilege}:
loading BPF LSM programs requires \texttt{CAP\_BPF} and installs global
policy, whereas envelopes attach unprivileged and per-launch, in the
spirit of Landlock's self-restriction model.  Second,
\emph{structural non-refreshability}: a BPF-map budget can be
re-incremented by any component that can reach the map---the verifier
guarantees memory safety, not that a policy never tops a counter back
up.  Envelope counters live in kernel-owned state, and the only
mutation exposed to user space is decrement-by-charge
(\S\ref{sec:design}).  Third, \emph{io\_uring}: fixed-file
registrations bypass the file-based LSM hooks, and submission-point
hooks cannot charge operations that execute later in kernel context;
our mediation sits at the SQE execution point and fails closed on
fixed files (\S\ref{sec:impl}).  A privileged BPF-LSM reimplementation
could approximate much of the envelope policy surface; we view such a
Linux port as worthwhile future work and position envelopes as the
unprivileged, deployment-time formulation with mechanically verified
coverage.

\paragraph{Sandbox runtimes.}
gVisor~\cite{gvisor2019} and WebAssembly~\cite{haas2017wasm} confine
unmodified (resp.~recompiled) code behind an interposed runtime.
gVisor demonstrates that user-space ABI interception is deployable at
the cost of a second system-call surface and non-trivial compatibility
and performance overheads; neither substrate provides budget, expiry,
or transfer-clamping dimensions.  Envelopes mediate in-kernel at
object granularity, so one policy engine covers both the Linux-ABI
compatibility path and native capability-aware services.

%%------------------------------------------------------------------------
\section{Design}
\label{sec:design}

\subsection{Dual-ABI Architecture}

A20OS exposes two system-call interfaces sharing a common kernel core:

\begin{itemize}
\item \textbf{Linux ABI} (361 entries): binary-compatible with musl,
      enabling unmodified static binaries to execute normally.
\item \textbf{Native ABI} (136 entries): capability-aware interface
      exposing handles, channels, and events for trusted services.
\end{itemize}

Both ABIs share the same VFS, memory manager, and scheduler.  The key
architectural property is that \emph{every} resource acquisition passes
through a common mediation layer regardless of which ABI issued the
request.  A capability envelope attaches to this mediation layer,
applying budget policies transparently to whichever ABI the enveloped
process uses.

\subsection{Budget Lattice}

Each envelope carries a policy structure defining per-class caps:

\begin{verbatim}
struct env_policy {
    uint32_t allowed_types;       /* class bitmask */
    uint64_t rights_cap[N];      /* per-type rights mask */
    uint64_t time_budget_ns;     /* temporal deadline */
    uint64_t op_budget;          /* total ops */
    uint64_t data_budget;        /* total bytes */
    uint32_t propagation_types;  /* outbound delegation */
};
\end{verbatim}

All budget counters decrease monotonically through mediated operations.
The mediator exposes no increment path: once set at creation time,
budgets can only be consumed, never refreshed.  This property---which
we call \emph{non-refreshability}---is enforced structurally: the only
mutation on budget counters is decrement-by-charge, and no syscall
increases any counter.

\subsection{Entry-Point Taxonomy}

We enumerate every path through which a process can acquire a new
usable resource authority (Table~\ref{tab:entrypoints}).  Each entry
is assigned a unique identifier (A1--A10) and mapped to a specific
kernel hook point where budget mediation fires.

\begin{table}[t]
\centering
\caption{Escape-surface entry points.}
\label{tab:entrypoints}
\small
\begin{tabular}{llp{3.2cm}}
\toprule
\textbf{ID} & \textbf{Path} & \textbf{Mediation} \\
\midrule
A1 & open/openat & Type $\cap$ rights $\cap$ cap \\
A2 & socket & Network rights + cap \\
A3 & pipe/memfd/eventfd/timerfd/signalfd & Creation-time \\
A4 & mmap (file-backed) & Map right + time \\
A5 & shmat & MEMORY class check \\
A6 & SCM\_RIGHTS recv/send & Transfer clamp / propagation \\
A7 & pidfd\_getfd & Fresh acquisition \\
A8 & /proc/self/fd reopen & Rights intersection \\
A9 & io\_uring fixed-file & Fail-closed \\
A10 & io\_uring SQE exec & Execution-point charge \\
\bottomrule
\end{tabular}
\end{table}

Coverage is mechanically verified: a generator script parses the
syscall table (365 entries), classifies each against the contract, and
a drift gate fails whenever a new syscall is registered without
explicit classification.  Currently 20 entries are actively mediated,
46 are classified PLANNED-W2 (scheduled refinement), and 299 involve
no resource authority.

\subsection{Direction Enforcement}

Each shadow handle tracks its granted rights as independent boolean
fields (read, write, stat, seek).  When an enveloped task performs an
I/O operation, the mediator checks the operation's direction against
the shadow's rights:

\begin{itemize}
\item READ-direction operations require the read bit;
\item WRITE-direction operations require the write bit;
\item Control operations (bind, listen, ioctl) skip direction checks.
\end{itemize}

This prevents a classic upgrade attack: opening a file read-only then
reopening it through \texttt{/proc/self/fd/N} with write flags.  The
reopen inherits only the intersection of requested and original rights,
so the write bit cannot be fabricated if absent from the source handle.

\subsection{Transfer Clamping}

When an enveloped process receives a descriptor via
\texttt{SCM\_RIGHTS} or \texttt{pidfd\_getfd}, the received authority's
rights are clamped to the intersection of the request and the
receiver's policy cap:

\[
  \text{granted} = \text{request} \cap \text{cap}
\]

An empty grant denies the transfer entirely.  This ensures transferred
authority never exceeds the receiver's policy ceiling, preventing both
upgrade (beyond policy) and fabrication (beyond request).

%%------------------------------------------------------------------------
\section{Implementation}
\label{sec:impl}

The mediator is implemented as a $\sim$670-line C module
(\texttt{kernel/ipc/envelope.c}) integrated into A20OS's system-call
dispatch layer.  It maintains per-envelope state including the policy
snapshot, live budget counters, and a shadow table keyed by global fd.

Control-plane syscalls (902--905) allow a supervisor to create
envelopes, attach processes, trigger revocation, and query audit
counters.  The deployment pattern follows a create-fork-enter-execve
model: the supervisor creates an envelope from a policy struct, forks,
and the child enters before calling \texttt{execve} on the untrusted
binary.  Once entered, the attachment is monotone: there is no leave
operation, and \texttt{execve} preserves the attachment.

USE-side hooks fire at six I/O families (read/write/readv/writev/
pread64/pwrite64) with per-family direction bits and conservative
data-budget pre-charging.  Vectorized variants pre-compute total
segment lengths before mediation so denial fires before any partial
transfer.

Socket operations mediate at three levels: creation (acquisition),
control (bind/connect/listen ops charge), and data-plane (send/recv
through separate hooks).  Accepted connections inherit the listener's
class but receive their own shadow entry clamped to the receiver's
policy.

io\_uring operations are mediated at the SQE execution point rather
than the submission point, ensuring that asynchronous operations
cannot bypass budget accounting after submission.  Fixed-file
registration is denied outright for enveloped tasks until a complete
transfer-tracking design is validated.

A coverage matrix generator mechanically parses the syscall table and
classifies each of 365 registered entries against the mediation
contract.  A drift gate fails compilation whenever a new syscall is
registered without explicit classification, preventing silent gaps in
the choke-point guarantee.

%%------------------------------------------------------------------------
\section{Evaluation}
\label{sec:eval}

We evaluate capability envelopes along three axes: security efficacy
(can attacks be blocked?), benign compatibility (do legitimate
workloads still function?), and overhead (what is the mediation cost?).

All experiments run under QEMU system emulation (riscv64, 1~GiB RAM)
on a host with AMD Ryzen 7735HS.  We deliberately report this fixed,
publicly reproducible configuration rather than hardware-specific
numbers: neither the security conclusions nor the relative-overhead
argument depends on the host, and emulator runs are exactly
reproducible by others.

\subsection{Attack Suite}

We exercise 15 scenarios---12 attack scenarios each verifying a specific
enforcement property, two benign multi-process flows, and one runtime
invariant audit (Table~\ref{tab:attacks}).  Every scenario runs
in a forked worker inside the envelope; the parent verifies the
expected outcome.

\begin{table}[t]
\centering
\caption{Attack suite results (all pass).}
\label{tab:attacks}
\small
\begin{tabular}{lp{4.2cm}l}
\toprule
\textbf{Scenario} & \textbf{Tests} & \textbf{Result} \\
\midrule
Type denial & socket() outside allowed\_types & EPERM \\
Rights deny & O\_WRONLY vs read-only cap & EACCES \\
Op budget & Third op after budget=2 exhausted & EACCES \\
Data budget & Write exceeding data\_budget & EACCES \\
Expiry & Post-deadline read after lazy sweep & EACCES \\
Revoke+kill & Active revoke + KILL\_ON\_EXPIRE & SIGKILL \\
Reopen upgrade & /proc/self/fd/N write-after-RD-only & EACCES \\
SCM receive & Descriptor from non-allowed class & dropped \\
SCM send & propagation\_types=0 sendmsg & EPERM \\
pidfd getfd & Cross-task fd theft & EPERM \\
shmat class & MEMORY outside allowed\_types & EPERM \\
socket data plane & Over-budget sendto on AF\_UNIX pair & EACCES \\
mksh flow & Full script under envelope & completed \\
pipe flow & 3-process pipeline, one envelope & completed \\
Audit walk & Post-exercise invariant check & zero viol. \\
\bottomrule
\end{tabular}
\end{table}

Every scenario passes: the mediator blocks all tested attacks while
permitting benign operations within policy bounds.  We emphasize that
these are mechanism tests we authored: they establish that each
enforcement property fires as designed, not adversarial completeness
against an independent attacker; \S\ref{sec:corpus} evaluates against
real malicious packages.

\subsection{Pilot Experiment: Coarse-Grained Dilemma}

We construct a five-scenario experiment comparing four defense arms
(Table~\ref{tab:pilot}).  Each cell forks a worker whose exit code
encodes the observed outcome.

\begin{table}[t]
\centering
\caption{Pilot experiment: five scenarios $\times$ four arms.}
\label{tab:pilot}
\small
\begin{tabular}{lp{1.2cm}p{1.2cm}p{1.2cm}p{1.2cm}}
\toprule
\textbf{Scenario} & \textbf{NONE} & \textbf{LL-perm} & \textbf{LL-strict} & \textbf{ENV} \\
\midrule
S0 install & ok & ok & \textbf{blocked} & ok \\
A1 exfil   & ok & ok & blocked & \textbf{blocked} \\
A2 loop    & ran & ran & ran & \textbf{stopped} \\
A3 escape  & ok & ok & blocked & pass$^*$ \\
A4 stall   & kept & kept & kept & \textbf{denied} \\
\bottomrule
\multicolumn{5}{l}{\footnotesize *A3: path dimension outside envelope v1 scope;}\\
\multicolumn{5}{l}{\footnotesize composes with Landlock for full coverage.}\\
\end{tabular}
\end{table}

The results quantify the coarse-grained dilemma: LL-permissive admits
all attacks (S0/A1/A2/A3/A4 all succeed); LL-strict blocks attacks
but also kills the benign install (S0 blocked).  Only the envelope
achieves both goals simultaneously: S0 completes while A1/A2/A4 are
blocked.  A3 (path escape) passes through the envelope v1 honestly
(no path dimension), composing with Landlock for combined coverage.
Each cell is a single deterministic run whose exit code encodes a
permit/deny outcome that is deterministic by construction; we
deliberately scope the claim to these five scenarios and do not
extrapolate to package-installation workloads in general.

\subsection{Corpus Evaluation: Real Malicious Packages}
\label{sec:corpus}

To move beyond self-authored scenarios, we evaluate against \emph{real}
malicious packages in two complementary ways: exact execution under
stock interpreters, and canonical replay for broader coverage.

\paragraph{Exact execution.}
From the DataDog malicious-software-packages dataset~\cite{datadogdataset}
we drew a seeded random sample of 100 packages (60 npm, 40 PyPI, from
pools of 12{,}601 and 1{,}816).  67 samples carry an install-time entry
point (34 npm lifecycle scripts, 33 PyPI \texttt{setup.py}); 66 executed
to completion in our matrix.  Samples run \emph{unmodified} on stock
Alpine Node.js 24.18.1 and CPython 3.12.14 inside A20OS (the
\texttt{pynode} instance), each twice: bare (NONE) and wrapped by an
install-time envelope (ENV: FILE+PIPE+event-queue+timer classes, no
SOCKET, 100k-op / 256-MiB budgets) via a \texttt{create}-\texttt{fork}-
\texttt{enter}-\texttt{execve} supervisor that also reads back the
mediator's audit counters.

Of the 66 samples, 25 reached a network operation at runtime---the
mediator's \texttt{acquire\_deny\_type} counter proves a
\texttt{socket()} call occurred---and \emph{all 25 were blocked} under
the envelope (Table~\ref{tab:exec}).  Under NONE, the same payloads
progressed to live network behavior: an exfiltration POST to a
\texttt{oastify.com} listener that sent data before the connection
broke, a completed TLS handshake to \texttt{webhook.site}, and DNS
resolution of an AWS Lambda C2 hostname; under ENV, each died at
\texttt{socket()} with EPERM.  The remaining 41 samples never reached
the network at runtime (payloads dormant, or crashing on interpreter
version/environment differences in \emph{both} arms)---a normal hit
rate for detonation of aged samples~\cite{duan2021measuring}.  No
interpreter aborted under the envelope, and no benign behavior was
denied.

\begin{table}[t]
\centering
\caption{Exact execution of 66 real malicious packages, NONE vs.\
ENVELOPE.  ``Reached net'' counts samples whose payload issued a
\texttt{socket()} (mediator counter).}
\label{tab:exec}
\small
\begin{tabular}{lcc}
\toprule
 & \textbf{NONE} & \textbf{ENVELOPE} \\
\midrule
Reached network, attack proceeds & 25 & 0 \\
Reached network, blocked at \texttt{socket()} & 0 & 25 \\
No network attempt at runtime & 41 & 41 \\
Interpreter aborted by policy & --- & 0 \\
\bottomrule
\end{tabular}
\end{table}

\paragraph{Canonical replay for coverage.}
Separately, a static classifier
(\texttt{tools/corpus/extract\_corpus.py}) maps 87 of the 100 samples to
a canonical behavior class (6 credential exfiltration, 25 network
beaconing, 56 second-stage download) without executing them.  Replaying
each class's syscall sequence under the same two arms confirms the
mechanism at scale: every attack goal succeeds under NONE (87/87,
validating replay fidelity) and is blocked under the envelope (87/87),
stable across three repetitions.  The install footprints of 20 popular
benign npm packages (120 files, 258\,KiB of real tarball writes)
complete under manifest-sized budgets with $2\times$ headroom (20/20).

We stress the boundary of each method: replay covers more samples but
executes canonical sequences, not the original code; exact execution
runs the original code but only as far as aged, partially dormant
payloads go in a fresh sandbox.  The two agree on the claim that
matters: whenever install-time code reaches for the network, the
envelope denies it.

\subsection{Real-Binary Compatibility}

The G1 mksh-flow cell proves that a real static musl binary
(/bin/mksh, ${\sim}$400KB) runs correctly inside an envelope across
an execve boundary.  The shell's internal openat/write operations are
mediated and tracked, and the marker file is written correctly.

The G2 wget-blocked cell demonstrates TYPE-level enforcement blocking
a real binary's network access.  Under an envelope whose
\texttt{allowed\_types} excludes SOCKET, \texttt{wget} exits nonzero
because its internal \texttt{socket()} call returns EPERM.

The G4 pipe-flow cell extends this to multi-process composition: an
enveloped shell builds a three-stage pipeline using builtins only.
Every segment inherits the shared envelope across fork, both
\texttt{pipe2()} acquisitions are mediated (A3), and the final
redirect's openat/write is charged -- three processes, one envelope,
zero unmediated steps.

\subsection{Overhead}

We measure per-operation latency with and without envelope activation
using CLOCK\_MONOTONIC sampled once per 20{,}000-iteration loop
(Table~\ref{tab:overhead}).

\begin{table}[t]
\centering
\caption{Mediation overhead (QEMU/TCG, directional).}
\label{tab:overhead}
\small
\begin{tabular}{lrrr}
\toprule
\textbf{Operation} & \textbf{off (ns)} & \textbf{on (ns)} & \textbf{$\Delta$} \\
\midrule
open+close & 79{,}587 & 82{,}834 & +4.1\% \\
read 64B   & 12{,}238 & 15{,}815 & +29.2\% \\
write 64B  & 32{,}049 & 37{,}731 & +17.7\% \\
lseek (ctl)& 11{,}383 & 11{,}507 & +1.1\% \\
sendto+drain 64B & 96{,}051 & 104{,}795 & +9.1\% \\
\bottomrule
\end{tabular}
\end{table}

The lseek control group shows noise-level variation (+1.1\%),
validating measurement methodology.  Acquisition-path overhead is
+4.1\%; file use-path overhead ranges from +17.7\% (write) to
+29.2\% (read).  The socket data plane, measured as a paired
\texttt{sendto} plus recvfrom drain over an AF\_UNIX stream pair,
adds +9.1\%: transport cost dominates the mediated pair, so relative
mediation overhead is lower than on the file path.  Absolute deltas
are ${\sim}3.6\mu$s (read) and ${\sim}5.7\mu$s (write) under TCG
emulation; all comparisons are same-configuration.

%%------------------------------------------------------------------------
\section{Formal Verification}
\label{sec:formal}

We verify the mediator's safety properties using two complementary
approaches.

\subsection{TLA+ Model Checking}

We encode the mediator state machine in TLA+
(\texttt{Envelope.tla}, 199 lines) and exhaustively check it with TLC
over a bounded state space (OpsMax=3, DataMax=4, ExpireAt=3, MaxTick=6,
3 tasks, 2 gfds, 2 types).  Five invariants hold:

\begin{itemize}
\item \textsc{TypeAllowed}: shadow types within allowed\_types;
\item \textsc{RightsSubCap}: shadow rights within class caps;
\item \textsc{OpsNonNeg} / \textsc{DataNonNeg}: budget counters bounded;
\item \textsc{ClockBounded}: clock within model horizon.
\end{itemize}

Two temporal properties hold: NoResurrect (expired stays expired) and
KillOnce (kill flag is one-way).  TLC explores 27{,}829 states
generating 4{,}620 distinct states at depth 14 in under one second.

\subsection{Lean 4 Proofs}

We formalize core budget lattice properties in Lean~4
(\texttt{BudgetLattice.lean}, 200 lines) using only core constructs
(no Mathlib dependency).  Thirteen theorems are proved with zero
\texttt{sorry}, in four families:

\begin{itemize}
\item \emph{Budget decay and expiry}: \texttt{deduct\_le}
      (monotone decay), \texttt{deduct\_strict} (strict positive
      decay), \texttt{decay\_transitive} (chained spending),
      \texttt{latch\_true} (one-way expiry latch);
\item \emph{Transfer clamping}:
      \texttt{and\_field\_le\_cap} (clamped grant never exceeds the
      policy cap), \texttt{and\_field\_preserves} (legitimate
      authority survives clamping), \texttt{and\_field\_idem}
      (idempotent under repeated mediation);
      \texttt{xferred\_wr/rd\_needs\_both} enforce direction on
      transferred shadows -- a write- or read-pass implies both
      requester and policy granted the bit;
\item \emph{Delegation bounding}:
      \texttt{del\_wr/rd\_needs\_both} -- exercising a delegated
      right requires delegator grant AND receiver cap, so authority
      cannot escalate through the trust chain;
\item \emph{Cross-envelope budgets}:
      \texttt{scm\_recv\_le\_receiver},
      \texttt{scm\_recv\_le\_sender} -- receive-time charging must
      stay within $\min(\mathrm{sender},\mathrm{recv})$; our v1
      prototype satisfies both bounds by construction (receipt draws
      one op from the receiver's own remaining budget and couples to
      no sender-side amount), so an incoming gift overdraws neither
      party.
\end{itemize}

\subsection{From Models to Code}
\label{sec:trace}

We do not claim refinement between the models and the C implementation.
Instead, each model-checked property is tied to a single auditable code
point in \texttt{kernel/ipc/envelope.c}, and---unusually---the same
invariants are \emph{re-checked at runtime}: the \texttt{env\_audit}
syscall walks every live envelope and re-verifies
\textsc{TypeAllowed}, \textsc{RightsSubCap}, and the budget bounds on
live kernel state (the \emph{Audit walk} row of
Table~\ref{tab:attacks}).  The correspondence is:

\begin{itemize}
\item \textsc{TypeAllowed}: the \texttt{allowed\_types} gate in
      \texttt{env\_mediate\_acquire} and \texttt{env\_mediate\_class}
      precedes shadow installation;
\item \textsc{RightsSubCap}: \texttt{env\_acquire\_charge\_install\_locked}
      stores \texttt{rights \& cap}, and the transfer path
      (\texttt{env\_mediate\_acquire\_gfd}) stores
      \texttt{want \& rights\_by\_class[kind]};
\item \textsc{OpsNonNeg}/\textsc{DataNonNeg}: every charge site checks
      the counter for zero (resp.~sufficiency) \emph{before}
      decrementing, so counters cannot underflow;
\item \textsc{NoResurrect}: the \texttt{expired} flag is set only by
      \texttt{env\_mark\_expired\_locked} and has no clear path;
\item \textsc{KillOnce}: \texttt{kill\_sent} is an atomic test-and-set
      shared by the expiry path and \texttt{env\_revoke};
\item Monotone attachment: \texttt{env\_enter} is a compare-and-swap
      from \texttt{NULL} only, so an enveloped task can neither switch
      envelopes nor leave.
\end{itemize}

The runtime audit turns the model-checked invariants into executable
checks on real kernel state, catching divergence between the models and
the implementation as a testable violation rather than a silent gap.

%%------------------------------------------------------------------------
\section{Related Work}
\label{sec:related}

\paragraph{Capability systems.}
KeyKOS~\cite{hardy1985keykos} and EROS~\cite{shapiro1999eros}
established pure capability operating systems with persistent,
delegatable authority, but required applications written against their
object model.  Capsicum~\cite{watson2010capsicum} brought capability
mode to a commodity UNIX at the cost of the rewrite mandate we target.
Zircon~\cite{zircon2023} exposes handle-based objects in a production
microkernel---our native ABI follows a similar shape---but its security
story targets new userspace, not legacy Linux binaries.
CHERI~\cite{woodruff2014cheri} provides hardware capabilities with
architectural support; envelopes are software-only and operate on
unmodified binaries.

\paragraph{Unprivileged self-restriction.}
OpenBSD's pledge~\cite{pledge2} narrows the syscall surface by promise
group but requires source modification and offers no budget dimensions;
unveil~\cite{unveil2} restricts path visibility but, like Landlock,
cannot express volume or time.
Landlock~\cite{landlock2023} provides unprivileged filesystem access
control on Linux; our pilot experiment quantifies its coarse-grained
dilemma.  seccomp-BPF~\cite{seccomp2012} filters by syscall number and
scalar arguments but cannot inspect pointer-indirect arguments and
carries no per-object state.

\paragraph{Mandatory access control and the LSM framework.}
The LSM framework~\cite{wright2002lsm} made Linux security policy
pluggable; SELinux~\cite{loscocco2001selinux} and
AppArmor~\cite{cowan2000subdomain} provide mandatory access control
behind it, requiring administrator-installed global policy and, in the
path-based AppArmor case, inheriting the path granularity we avoid.
BPF LSM~\cite{krsi2019} adds programmability at the same hooks;
\S\ref{sec:positioning} details the structural differences from
envelopes (deployment privilege, non-refreshability, io\_uring).

\paragraph{Resource accounting.}
Resource containers~\cite{banga1999rc} introduced kernel-visible
resource principals for consumption accounting; POSIX
rlimits~\cite{setrlimit2017} and cgroups v2~\cite{cgroupsv2} are their
deployed descendants.  All three meter consumption rather than
authority and provide neither direction enforcement nor transfer
clamping (\S\ref{sec:positioning}).

\paragraph{Sandbox runtimes.}
Native Client~\cite{yee2009nacl} and WebAssembly~\cite{haas2017wasm}
confine code through software fault isolation at the cost of
recompilation or restricted toolchains.  gVisor~\cite{gvisor2019}
interposes a user-space kernel for unmodified containers, trading a
second system-call surface and compatibility overhead for isolation;
it has no budget or expiry dimensions.

\paragraph{Supply-chain security.}
Measurement studies document the scale and install-time character of
open-source supply-chain attacks~\cite{ohm2020backstabber,duan2021measuring,zimmermann2019smallworld,ladisa2023taxonomy}.
On the defense side, Latch~\cite{latch2022} infers install-time
manifests via strace and enforces them via AppArmor, and
Leash~\cite{leash2025} sandboxes installer scripts using FreeBSD
Capsicum with a user-space supervisor.  Both operate at coarser
granularity than object-level mediation and lack budget dimensions.

\paragraph{Formal verification of OS kernels.}
seL4~\cite{klein2009sel4} achieved full functional-correctness
verification of a microkernel; CertiKOS~\cite{gu2016certikos} verified
concurrent kernels via layered refinement; Serval~\cite{nelson2019serval}
automated verification of systems code through symbolic evaluation.  We
target specific security properties of one subsystem rather than
whole-kernel refinement, using TLA+ model
checking~\cite{lamport2002specifying} and Lean~4
proofs~\cite{demoura2021lean4} as complementary tools, with runtime
re-verification bridging models and code (\S\ref{sec:trace}).

%%------------------------------------------------------------------------
\section{Discussion and Limitations}
\label{sec:discussion}

\paragraph{Scope.}
Envelopes constrain \emph{authority lifetime and volume}, not
information-flow control.  Data willingly written to permitted
destinations persists beyond envelope expiry.  This matches the threat
model: the mechanism bounds unauthorized acquisition, operation
volume, and time window---not post-hoc recall.

\paragraph{Path dimension.}
Envelope v1 mediates by resource class, not by filesystem path.
Path-scoped policies compose orthogonally with envelopes (e.g.,
Landlock for WHERE, envelope for HOW-MUCH/HOW-LONG).

\paragraph{Performance.}
Overhead figures come from the fixed QEMU/TCG configuration above;
TCG amplifies syscall-dense loops, so absolute values should be read
as same-configuration comparisons rather than native-hardware costs.
Application-level end-to-end benchmarks remain future work.

\paragraph{Evaluation breadth.}
The corpus evaluation (\S\ref{sec:corpus}) grounds the security claim
in real malicious packages executed on stock interpreters, but two
boundaries remain.  First, aged samples detonate imperfectly: 41 of 66
payloads never reached the network in \emph{either} arm (dormancy,
interpreter drift, dead infrastructure), so our block rate is measured
on the 25 samples that demonstrably attacked.  Second, all 25 were
network-dependent; path-only classes (persistence writes, credential
reads without exfiltration) are outside envelope v1's scope and compose
with Landlock, as in the pilot's A3.  A privileged BPF-LSM port of the
envelope policy surface would enable same-host comparison against
commodity Linux (\S\ref{sec:positioning}).

\paragraph{PLANNED-W2 surfaces.}
Forty-six syscalls are classified PLANNED-W2 (unmediated but
enumerated).  Key items include socket data-plane charging
(sendto/recvfrom), mmap file-backed mapping, and SysV IPC creation.
Each is tracked in the coverage matrix with assigned W2 line items.

%%------------------------------------------------------------------------
\section{Conclusion}
\label{sec:conclusion}

We presented capability envelopes: a kernel-level mechanism that
attaches quantified budget policies to unmodified Linux binaries at
deployment time.  Unlike prior capability systems that require source
modification, envelopes exploit the asymmetry that budgets are
deployable even when permissions are not.  A dual-ABI kernel
architecture enables transparent enforcement across 366 registered
syscalls with direction-aware right checks, conservative data-budget
pre-charging, and mechanical coverage verification.

Our pilot evaluation demonstrates that envelopes resolve the
coarse-grained dilemma that plagues path-based policies \emph{within
the tested scenarios}: the envelope arm is the only configuration that
both completes the benign installation and blocks the tested attack
classes, through type denial, rights clamping, and budget exhaustion.
Exact execution of 66 real malicious packages on stock interpreters
confirms the result at corpus scale: every sample that reached the
network (25/25) was denied at \texttt{socket()}, while undefended
payloads reached live C2 infrastructure; an 87-sample canonical replay
shows the same boundary.
Overhead under emulation is $+4.1\%$ on the acquisition path and
$+17.7$--$29.2\%$ on the file use path.  A TLA+ model exhaustively
verifies safety over a bounded state space, Lean proofs establish core
lattice properties without external dependencies, and a runtime audit
re-checks the model-checked invariants on live kernel state.

\paragraph{Availability.}
The A20OS kernel, the envelope mediator, the attack-suite and pilot
harnesses, the TLA+ model, and the Lean~4 proofs are open source and
publicly reproducible under the fixed QEMU configuration reported in
\S\ref{sec:eval}.

%%------------------------------------------------------------------------
\begin{thebibliography}{10}

\bibitem{watson2010capsicum}
R.~N.~M. Watson, J. Anderson, B. Laurie, and K. Kennaway,
``Capsicum: Practical capabilities for UNIX,''
in \emph{Proc.\ 19th USENIX Security Symposium}, 2010.

\bibitem{hardy1985keykos}
N. Hardy,
``The KeyKOS architecture,''
\emph{ACM Operating Systems Review}, 19(4):8--25, 1985.

\bibitem{shapiro1999eros}
J.~S. Shapiro, J.~M. Smith, and D.~J. Farber,
``EROS: A fast capability system,''
in \emph{Proc.\ 17th ACM Symposium on Operating Systems Principles (SOSP)}, 1999.

\bibitem{zircon2023}
Google,
``Fuchsia: Zircon kernel and handle-based object model,''
Fuchsia project documentation, 2023. \url{https://fuchsia.dev/fuchsia-src/concepts/kernel}.

\bibitem{woodruff2014cheri}
J. Woodruff, R.~N.~M. Watson, D. Chisnall, S.~W. Moore, J. Anderson,
B. Davis, B. Laurie, P.~G. Neumann, R. Norton, and M. Roe,
``The CHERI capability model: Revisiting RISC in an age of risk,''
in \emph{Proc.\ 41st International Symposium on Computer Architecture (ISCA)}, 2014.

\bibitem{pledge2}
OpenBSD,
``\texttt{pledge(2)} --- restrict system operations,''
OpenBSD Programmer's Manual, 2015.

\bibitem{unveil2}
OpenBSD,
``\texttt{unveil(2)} --- unveil parts of a directory tree,''
OpenBSD Programmer's Manual, 2018.

\bibitem{landlock2023}
M. Sala\"un,
``Landlock: Unprivileged access control,''
Linux kernel documentation, v6.2+, 2023.

\bibitem{seccomp2012}
W. Drewry,
``Seccomp BPF (SECure COMPuting with filters),''
Linux kernel documentation, v3.5+, 2012.

\bibitem{wright2002lsm}
C. Wright, C. Cowan, S. Smalley, J. Morris, and G. Kroah-Hartman,
``Linux Security Modules: General security support for the Linux kernel,''
in \emph{Proc.\ 11th USENIX Security Symposium}, 2002.

\bibitem{loscocco2001selinux}
P. Loscocco and S. Smalley,
``Integrating flexible support for security policies into the Linux operating system,''
in \emph{Proc.\ USENIX Annual Technical Conference (Freenix Track)}, 2001.

\bibitem{cowan2000subdomain}
C. Cowan, S. Beattie, G. Kroah-Hartman, C. Pu, P. Wagle, and V. Gligor,
``SubDomain: Parsimonious server security,''
in \emph{Proc.\ 14th USENIX Conference on System Administration (LISA)}, 2000.

\bibitem{krsi2019}
K.~P. Singh,
``Kernel Runtime Security Instrumentation (BPF LSM),''
Linux kernel documentation, 2019.

\bibitem{banga1999rc}
G. Banga, P. Druschel, and J.~C. Mogul,
``Resource containers: A new facility for resource management in server systems,''
in \emph{Proc.\ 3rd USENIX Symposium on Operating Systems Design and Implementation (OSDI)}, 1999.

\bibitem{setrlimit2017}
IEEE and The Open Group,
``\texttt{setrlimit()} --- control maximum resource consumption,''
POSIX.1-2017, 2017.

\bibitem{cgroupsv2}
The Linux kernel documentation,
``Control Group v2,''
online. \url{https://docs.kernel.org/admin-guide/cgroup-v2.html}.

\bibitem{yee2009nacl}
B. Yee, D. Sehr, G. Dardyk, J.~B. Chen, R. Muth, T. Ormandy,
S. Okasaka, N. Narula, and N. Fullagar,
``Native Client: A sandbox for portable, untrusted x86 native code,''
in \emph{Proc.\ 30th IEEE Symposium on Security and Privacy (S\&P)}, 2009.

\bibitem{gvisor2019}
Google,
``gVisor: A container sandbox runtime,''
online, 2019. \url{https://gvisor.dev/docs}.

\bibitem{haas2017wasm}
A. Haas, A. Rossberg, D.~L. Schuff, B.~L. Titzer, M. Holman,
D. Gohman, L. Wagner, A. Zakai, and J. Bastien,
``Bringing the Web up to speed with WebAssembly,''
in \emph{Proc.\ 38th ACM Conference on Programming Language Design and Implementation (PLDI)}, 2017.

\bibitem{ohm2020backstabber}
M. Ohm, H. Plate, A. Sykosch, and M. Meier,
``Backstabber's knife collection: A review of open source software supply chain attacks,''
in \emph{Proc.\ 17th International Conference on Detection of Intrusions and Malware, and Vulnerability Assessment (DIMVA)}, 2020.

\bibitem{duan2021measuring}
R. Duan, O. Alrawi, R.~P. Kasturi, R. Elder, B. Saltaformaggio, and W. Lee,
``Towards measuring supply chain attacks on package managers for interpreted languages,''
in \emph{Proc.\ 28th Network and Distributed System Security Symposium (NDSS)}, 2021.

\bibitem{zimmermann2019smallworld}
M. Zimmermann, C.-A. Staicu, C. Tenny, and M. Pradel,
``Small world with high risks: A study of security threats in the npm ecosystem,''
in \emph{Proc.\ 28th USENIX Security Symposium}, 2019.

\bibitem{ladisa2023taxonomy}
P. Ladisa, H. Plate, M. Martinez, and O. Barais,
``SoK: Taxonomy of attacks on open-source software supply chains,''
in \emph{Proc.\ 44th IEEE Symposium on Security and Privacy (S\&P)}, 2023.

\bibitem{datadogdataset}
DataDog,
``Malicious software packages dataset,''
online, 2023--2026. \url{https://github.com/DataDog/malicious-software-packages-dataset}.

\bibitem{latch2022}
Z. Luo et al.,
``Latch: Preventing install-time attacks in npm with lightweight manifest checking,''
in \emph{Proc.\ 17th ACM Asia Conference on Computer and Communications Security (AsiaCCS)}, 2022.

\bibitem{leash2025}
S. Jadidi and J. Anderson,
``Leash: A transparent capability-based sandboxing supervisor for Unix,''
in \emph{Proc.\ 11th International Conference on Information Systems Security and Privacy (ICISSP)}, 2025.

\bibitem{klein2009sel4}
G. Klein, K. Elphinstone, G. Heiser, J. Andronick, D. Cock, P. Derrin,
D. Elkaduwe, K. Engelhardt, R. Kolanski, M. Norrish, T. Sewell,
H. Tuch, and S. Winwood,
``seL4: Formal verification of an OS kernel,''
in \emph{Proc.\ 22nd ACM Symposium on Operating Systems Principles (SOSP)}, 2009.

\bibitem{gu2016certikos}
R. Gu, Z. Shao, H. Chen, X.~N. Wu, J. Kim, V. Sj\"oberg, and D. Costanzo,
``CertiKOS: An extensible architecture for building certified concurrent OS kernels,''
in \emph{Proc.\ 12th USENIX Symposium on Operating Systems Design and Implementation (OSDI)}, 2016.

\bibitem{nelson2019serval}
L. Nelson, J. Bornholt, R. Gu, A. Baumann, E. Torlak, and X. Wang,
``Scaling symbolic evaluation for automated verification of systems code with Serval,''
in \emph{Proc.\ 27th ACM Symposium on Operating Systems Principles (SOSP)}, 2019.

\bibitem{lamport2002specifying}
L. Lamport,
\emph{Specifying Systems: The TLA+ Language and Tools for Hardware and Software Engineers},
Addison-Wesley, 2002.

\bibitem{demoura2021lean4}
L. de Moura and S. Ullrich,
``The Lean 4 theorem prover and programming language,''
in \emph{Proc.\ 28th International Conference on Automated Deduction (CADE)}, 2021.

\end{thebibliography}

\end{document}
